\documentclass[11pt]{article}
\usepackage[margin=1in]{geometry}
\usepackage[T1]{fontenc}
\usepackage{amsmath,amssymb,array,hyperref}
\setlength{\emergencystretch}{2em}
\newcommand{\simplex}{\Delta_J}
\title{Temporally Regularized Synthetic Control in R:\\
A Chronological Workflow}
\author{}
\date{September 15, 2026}
\begin{document}
\maketitle
\begin{abstract}
This companion paper specifies a chronological workflow for a proposed R
implementation of temporally regularized synthetic control. It separates
panel preparation, predetermined predictor metrics, donor-weight fitting,
chronological validation, counterfactual prediction and treatment-effect
calculation. The planned interface separates panel preparation, optional
chronological split construction, fitting and effect estimation, with a
reusable prediction method. This initial design draft preserves the agreed
software specification while the companion methodological paper is being
settled. It does not report an implemented package, timing benchmarks or
empirical results. Software development and validation will follow the
methodological review; the California application will then be reported
in the methodological paper.
\end{abstract}

\section{Purpose and scope}
The companion methodological paper, \emph{Temporally Regularized Synthetic
Control}, develops the estimators, conditional theory and proofs, and will
report the empirical application. This paper has a different purpose:
to document the data interfaces, information restrictions, numerical
implementation and reproducible workflow that realize those estimators.
Neither paper should be read as evidence that the proposed software has
already been implemented or that counterfactual accuracy has been verified.

The development sequence is to settle the methodological specification,
then implement and validate the software described here, then use those
functions to complete the empirical comparison in the methodological
paper. The following material is a retained design specification, not a
parallel implementation task. Package naming, dependencies, empirical
predictor-weight calibration and numerical acceptance settings remain to
be settled before release. The brief mathematical summary makes the
workflow self-contained; the proofs remain in the methodological paper.

\section{Estimator interface}\label{sec:framework}
Consider one target and $J\geq2$ eligible donors observed at $T\geq3$
equally spaced pre-treatment dates. Let $X_{1,t}\in\mathbb R^K$ and
$X_{0,t}\in\mathbb R^{K\times J}$ be aligned predictor blocks, with
common predictor definitions and donor ordering. The researcher supplies
$V_t=\operatorname{diag}(v_{1,t},\ldots,v_{K,t})$, where $v_{k,t}\geq0$
and $\sum_kv_{k,t}=1$, before donor fitting. Uniform, empirical constant
and empirical time-varying metrics are distinct input modes, not variables
jointly optimized with donor weights.

Write $\simplex=\{w\in\mathbb R^J:w\geq0,\mathbf1^\top w=1\}$,
$W=(w_1,\ldots,w_T)\in\mathcal C=\simplex^T$, and
$\theta=\operatorname{col}_t(w_t)$. Define
$\Delta w_t=w_t-w_{t-1}$ and
$\Delta^2w_t=w_t-2w_{t-1}+w_{t-2}$. For $q\in\{1,2\}$,
\[
 A_{\mathcal V}(W)=T^{-1}\sum_{t=1}^T
       \|X_{1,t}-X_{0,t}w_t\|_{V_t}^2,\qquad
 C_q(W)=(T-q)^{-1}\sum_{t=q+1}^T\|\Delta^q w_t\|_2^2.
\]
The fitted path minimizes
\[
 F_{q,\lambda,\mathcal V}(W)=(1-\lambda)A_{\mathcal V}(W)
                      +\lambda C_q(W),\qquad 0\leq\lambda<1.
\]
Among primary minimizers, select the one of smallest squared Euclidean
norm. This secondary selection is not an additional ridge term.
The additive coefficient is $\alpha=\lambda/(1-\lambda)$;
$\lambda=1$ is excluded because it discards matching entirely.

For the numerical formulation below, define
\[
 \mathcal A\theta=\operatorname{col}_t
       (T^{-1/2}V_t^{1/2}X_{0,t}w_t),\qquad
 b=\operatorname{col}_t(T^{-1/2}V_t^{1/2}X_{1,t}),
\]
\[
 \mathcal D_q\theta=(T-q)^{-1/2}
       \operatorname{col}_{t=q+1}^T(\Delta^q w_t).
\]
Let $\widehat a=\widehat w_T$ and
$\widehat d=\widehat w_T-\widehat w_{T-1}$ for an accepted fitted path.
For either penalty order, define endpoint retention (E) and projected
increment continuation (D) by
\begin{equation}\label{eq:continuation}
 \widehat w^E_{T+h}=\widehat a,\qquad
 \widehat w^D_{T+h}=\Pi_{\simplex}(\widehat a+h\widehat d),
 \quad h=1,\ldots,H,
\end{equation}
where $\Pi_{\simplex}$ is Euclidean projection. E retains the final
composition rather than returning to a separately fitted static vector.
D continues the final increment and enforces feasibility at each horizon.
Projected increment continuation is the package default for both orders
and the continuation rule in the primary empirical comparison. Endpoint
retention remains an optional within-fit diagnostic.
D uses the original endpoint and increment, not recursively projected
updates.\footnote{Pointwise projection of an affine candidate is not in
general the solution to a joint minimum-curvature problem over the entire
future path. We specify it as a continuation rule, not as an automatic
consequence of the training optimization.}


Here $x_{T+h}$ denotes aligned contemporaneous donor outcomes in original
units, and the untreated prediction is
$\widehat Y^R_{1,T+h}(0)=x_{T+h}^\top\widehat w^R_{T+h}$ for
$R\in\{E,D\}$. The observed treated target outcome is
$Y^{\rm obs}_{1,T+h}$; $\tau_{T+h}=Y_{1,T+h}(1)-Y_{1,T+h}(0)$ is its
unobserved causal effect. In validation windows $y_t$ denotes the observed
untreated target outcome. Prediction does not require future target
predictors or target outcomes.

\section{Estimation Strategy}\label{sec:estimation}
This section gives the chronological workflow for the proposed R package.
The revised calibration and interface decisions below were agreed on
September 15, 2026. A first scoped refactor now integrates version-2 raw
blocks into preparation and updates split-input compatibility, with revised
preparation and split fixtures. The user-supplied September 15, 2026 R
console transcript shows both revised suites passing, including explicit
source loading and an isolated-environment preparation check. Session
information and source revision identifiers were not supplied, and the
transcript does not establish a clean-start session. Independent local
reproduction remains pending because R is unavailable in this workspace.
The separate \texttt{tvsc.vmetrics()} source draft now implements the
approved scaling and calibration contract. A subsequent user-supplied
excerpt shows its final isolated-loading check and suite success message;
the full fixture log and independent reproduction remain pending.
Fitting, prediction and effect operations remain unimplemented.
Earlier test reports describe the preceding interface.

\subsection{Revised interface and calibration design}\label{sec:revised-design}
The public preparation operation will combine panel validation with raw
block construction. The researcher declares target and donor identities;
\texttt{tvsc.dataprep()} validates them and arranges the corresponding
unscaled predictor and outcome arrays. Block construction may remain an
internal helper. Preparation does not learn scales or predictor relevance.
Post-treatment outcomes remain separate from all fitting inputs.

\paragraph{Approved preparation interface.}
Retain the existing argument names and the default regular step:
\begin{verbatim}
tvsc.dataprep(data, unit, time, outcome, treated, donors,
              intervention_time, pre_period, predictors,
              period_step = 1)
\end{verbatim}
The explicit \texttt{predictors} argument lists exactly the matching
features; the outcome is not added silently. For example, with
\texttt{outcome = "sales"}, \texttt{predictors = "sales"} prepares the
approved outcome-only specification, whereas
\texttt{predictors = c("sales", "income")} includes both features.
The intervention index denotes the first treated period, and
\texttt{pre\_period} explicitly lists regularly spaced training dates.
Retain supplied donor and predictor ordering and the existing training
validation rules.

Retain the \texttt{tvsc\_prepared} class and its top-level
\texttt{schema\_version}, \texttt{schema}, \texttt{training},
\texttt{future\_donors}, \texttt{future\_target}, \texttt{recipe} and
\texttt{diagnostics} components, adding \texttt{blocks}. This component
contains the unscaled \texttt{X1}, \texttt{X0}, \texttt{Y1} and
\texttt{Y0} arrays for the full designated pre-period, with dates and
dimensions. Preserve the existing block layout: date-indexed target
feature vectors and predictor-by-donor matrices, a target outcome vector,
and a date-by-donor outcome matrix. Keep the ordered raw training table
so subsequent operations can obtain prefixes and reconstruct preprocessing
rather than slice full-period scaled data. The recipe records
contemporaneous features and unapplied scaling; diagnostics retain
validation results and constant-feature flags without deleting features.
Available future outcomes remain separate, and future target outcomes
are optional. Preparation performs no scaling, calibration, splitting,
donor fitting or effect calculation. A separate public
\texttt{tvsc.blocks()} call is no longer required in the revised workflow.
These arguments and output responsibilities are now reflected in the
preparation source draft, with a user-run pass of its revised fixtures.
Independent reproduction and broader integration validation remain pending.
The versioning rule follows.

\paragraph{Approved object versions and compatibility.}
Version each object type separately. Revised \texttt{tvsc\_prepared}
objects, raw and scaled block objects, and scaling metadata use
\texttt{schema\_version = 2L}. The new \texttt{tvsc\_vmetrics} class
starts at \texttt{1L}. Split objects retain class \texttt{tvsc\_splits}
and version \texttt{1L}, because their date-based output structure does
not change. These numbers identify object contracts, not package releases
or evidence of validation.

Revised preparation produces version-2 objects. Metric construction
requires revised preparation and does not silently upgrade old objects.
Update the split constructor to accept revised preparation while keeping
its split-output format. Reject incompatible preparation objects with a
clear instruction to rerun preparation from the original panel; manually
changing a version field is not migration. Check required fields,
dimensions, ordering and relevant metadata in addition to class and
version. A version number alone does not establish object validity.

The first implementation step is limited to integrating raw block
construction into preparation, updating split-input compatibility, and
updating preparation and split tests. Metric construction follows as a
separate implementation step; the donor solver is outside this refactor.
Preparation now emits version-2 raw blocks and the split constructor
accepts version-2 preparation while retaining version-1 split output.
Their revised fixtures pass in the supplied user-run transcript; independent
reproduction remains pending. Version-2 scaling and
the new metric object are now constructed by the separate
\texttt{tvsc.vmetrics()} source draft, with a user-reported pass supported
by the final test excerpt. The old standalone block,
scaling, metric and problem drafts still expect version-1 inputs and are
not compatible with revised preparation; they are not silently migrated.

\paragraph{Metric construction and fitting responsibilities.}
The agreed name \texttt{tvsc.vmetrics()} denotes the revised operation
that combines scaling and predictor-metric construction for one specified
training sample. It returns both transformed blocks and their metrics,
with transformation parameters, calibration outputs and diagnostics.
Scaling is applied once, including when uniform metrics are requested.
The same transformed features are used for relevance calibration and
subsequent donor fitting; the calibration routine must not silently
standardize them again. The new source implements this replacement
separately; it does not rename or migrate the legacy \texttt{tvsc.metrics()}
interface.

\paragraph{Scoped metric-construction implementation.}
The source \texttt{R/tvsc.vmetrics.R} requires only the revised preparation
source, not the split constructor or legacy block/scaling functions.
It reconstructs a prefix of at least three dates from the raw training
table and checks agreement with its stored raw blocks. Full designated
schema metadata are retained, while periods, cutoff, arrays and learned
parameters describe only that prefix. Later measurement values and future
outcomes do not enter this operation. Internal helpers perform donor-only
scaling and SVD-based OLS; failures carry a stage and available diagnostics
in a \texttt{tvsc\_vmetrics\_error} condition. Named fallback records identify
the predictors and dates using divisor 1. The fixtures in
\texttt{tests/test-tvsc.vmetrics.R} cover known slopes and intercepts,
agreement with \texttt{lm.fit()}, both conversions, final-date metric
retention, donor-only and prefix isolation, supplied weights, numerical
failures and isolated source loading. On September 15, 2026 the user
supplied the suite ending: preparation and metric construction are sourced
into an environment with \texttt{baseenv()} as parent, and the displayed
assertion confirms agreement of constant-mode output with the existing
\texttt{scaled\_constant} reference, without local split or standalone
scaling bindings. The assertion completes without a displayed error and
is followed by the suite success message. This supports a user-reported
pass but does not establish execution of earlier assertions or creation
of the reference object in this run. The full log, session information,
source identifiers, clean-start evidence and independent reproduction
remain outstanding. The donor solver is unchanged.

\paragraph{Approved metric-construction interface.}
The approved core call is
\begin{verbatim}
tvsc.vmetrics(prepared, cutoff, predictor_weights = "uniform",
              scaling = "sd", conversion = "absolute",
              scales = NULL, provenance = NULL, rank_tol = NULL,
              weight_sum_tol = 1e-8)
\end{verbatim}
The mandatory \texttt{cutoff} is an actual date in the designated
pre-period, not a row number. Use only the training prefix ending there;
the final fit supplies the last pre-treatment date, and validation supplies
each fold's cutoff. No split object enters this operation.
The three character inputs for \texttt{predictor\_weights} are exactly
\texttt{"uniform"}, \texttt{"constant"} and \texttt{"varying"}.
They request equal feature weights, donor-only OLS with common slopes and
date intercepts, and date-by-date donor OLS with final-date metric
retention, respectively. The longer empirical-prefixed strings are not
part of the approved interface.

Alternatively, \texttt{predictor\_weights} accepts a named numeric vector
of supplied constant weights or a named predictor-by-date matrix of
supplied weights. Names must cover the selected features and, for a
matrix, exactly the prefix dates. Supplied weights refer to the selected
scaled feature coordinates; do not silently normalize them or slice extra
dates. The \texttt{conversion} choices are \texttt{"absolute"} and
\texttt{"squared"}, used only for empirical calibration and recorded as
not applicable otherwise. The \texttt{scales} input is required for
\texttt{scaling = "custom"} and must be \texttt{NULL} otherwise.
Outcome-only preparation uses \texttt{"uniform"}, returns $V_t=[1]$ and
records calibration bypass by design. Explicit \texttt{"constant"} or
\texttt{"varying"} requests for outcome-only inputs are rejected rather
than silently ignored.

Return an object of class \texttt{tvsc\_vmetrics} with
\texttt{schema\_version} and \texttt{schema} for feature definitions and
ordered unit identities, and \texttt{periods} and \texttt{cutoff} for the
actual training prefix. Its \texttt{blocks} contain transformed
\texttt{X1} and \texttt{X0} with unchanged training \texttt{Y1} and
\texttt{Y0}; its \texttt{weights} are a named $K\times T_c$ matrix whose
columns define $V_t$. The \texttt{scaling} component records the method,
centers, divisors, reference donors and fallback records. The
\texttt{calibration} component records mode, coefficients and intercepts,
calibration dates, conversion and boundary convention, or the reason
calibration was not performed. Retain \texttt{diagnostics} and
\texttt{provenance}; declared sources of external inputs do not certify
absence of leakage. The object contains no fitted donor weights, selected
penalty or validation scores and can be reused across penalty values and
both penalty orders for the same prefix and configuration. This core
interface is now reflected in the source draft. The rank and supplied-weight
sum controls are specified below, with object versions specified above.

\paragraph{Approved supplied-weight sum check.}
The default \texttt{weight\_sum\_tol = 1e-8} is an absolute tolerance.
For each supplied weight column require finite nonnegative entries, a
positive finite computed sum, and
\[
 \left|\sum_{k=1}^K v_{k,t}-1\right|
 \leq \texttt{weight\_sum\_tol}.
\]
Every column must pass; otherwise reject the supplied input. Even small
negative entries are rejected rather than clipped. Accepted values remain
unchanged, with no automatic renormalization. An override must be a finite
numeric scalar in $[0,1)$; zero requests exact equality of the computed
sum to one. Record the tolerance, column sums and maximum absolute
deviation from one in diagnostics. This is numerical acceptance, not
exact mathematical normalization: accepted supplied columns can differ
slightly from a sum of one. It is distinct from the intentional
normalization of OLS slopes when constructing empirical weights.
The source implements this check; R validation remains pending.

\paragraph{Approved numerical OLS rank check.}
For \texttt{"constant"}, subtract each calibration date's donor feature
means and stack the resulting donor rows across calibration dates. For
\texttt{"varying"}, form this centered donor design separately at each
calibration date. Write the resulting matrix as $B\in\mathbb R^{n\times K}$.
Centering accounts for the agreed date intercepts; it does not change the
matching blocks or introduce another feature standardization. Do not
rescale the design columns silently for this check.

Let $\sigma_r(B)$ be its singular values and let $\epsilon_{\rm machine}$
denote machine precision. The default \texttt{rank\_tol = NULL} selects
the relative tolerance $\rho=\max(n,K)\epsilon_{\rm machine}$. An override
must be a finite numeric scalar strictly between zero and one and replaces
$\rho$. Define
\[
 \eta=\rho\,\sigma_{\max}(B),\qquad
 r_{\rm num}=\#\{r:\sigma_r(B)>\eta\}.
\]
Require $r_{\rm num}=K$ for each required design; an all-zero design
has numerical rank zero and fails. Rejection causes explicit calibration
failure, not predictor deletion, reduced-rank coefficients, ridge or
uniform weights. Record the requested and actual tolerance, threshold,
singular values, numerical and required ranks, calibration dates and a
condition-number diagnostic. No additional condition-number rejection
threshold is introduced. Declare overrides before validation; do not
adjust them automatically until a failing fold passes. Passing this
roundoff-based check does not guarantee stable or substantively meaningful
predictor weights and can admit nearly dependent designs. This numerical
admissibility convention is not certification of exact mathematical rank.
The source implements this check; R validation remains pending.

\paragraph{Fitting orchestration.}
\texttt{tvsc.fit()} will construct chronological splits from explicit
window settings or accept a separately constructed \texttt{tvsc.split()}
object. The split helper remains available for inspection and reuse;
split construction does not belong inside \texttt{tvsc.vmetrics()}.
Within each fold, fitting obtains raw training-prefix blocks, reconstructs
scaling and metrics, freezes those metrics, then fits donor weights and
scores continuation predictions. Final fitting repeats these operations
on the full designated pre-period. A fixed penalty needs no validation
split. Integrating these calls does not introduce joint or alternating
optimization of predictor and donor weights. Exact revised fitting
arguments and return schemas remain to be specified.

\paragraph{Donor-only scaling and future-outcome calibration.}
Across uniform, empirical constant, empirical time-varying, supplied-weight
and outcome-only specifications, learned scaling parameters use donors
only within the actual training prefix. Apply the same transformation to
donors and target. No target-inclusive reference option is planned for
the first version. This common preprocessing rule separates the scaling
reference from predictor-weighting choices; it is not a claim of superior
counterfactual accuracy. The four scaling
methods remain \texttt{sd}, \texttt{zscore}, \texttt{none} and
\texttt{custom}, with \texttt{sd} the default. For learned within-date
sample variances the reference population changes from $J+1$ units to
$J$ donors, with denominator $J-1$. The \texttt{sd} method retains its
RMS within-date SD definition without centering; \texttt{zscore} retains
date-specific donor centers and scales. The \texttt{none} method leaves
features unchanged and estimates no reference parameters; \texttt{custom}
applies supplied positive scales identically to donors and target without
estimating them. Reconstruct learned preprocessing separately within each
validation training prefix and for the final full-pre-period fit.
Reporting outcomes are unchanged.
The legacy standalone scaling source still uses target plus donors; the
new internal scaling helper uses donors only and awaits R validation.
In particular, equal donor values do not imply a
zero target--donor contrast; a zero-scale rule must not erase that gap.

\paragraph{Approved zero-scale convention.}
Use divisor 1 for genuine zero donor variation and record a visible
diagnostic identifying the affected features and dates. For \texttt{sd},
this fallback applies only when the feature is constant across donors at
every training date; the common donor value may differ across dates.
A zero-variation date alone does not trigger fallback when the common RMS
scale is positive. No centering is introduced. For \texttt{zscore}, apply
the rule separately to each feature/date, retaining its donor center and
using divisor 1. Apply the same transformation to target and donors:
equal donor values must not force a different target value to zero.
Affected matching gaps remain in original feature units, not unit-variance
units. The \texttt{none} method remains unchanged; \texttt{custom} requires
strictly positive finite supplied scales and rejects zero rather than
replacing it.

The fallback requires actual donor equality, not a small computed SD.
Unequal donor values producing a zero or nonfinite scale are a numerical
failure, not a reason to apply divisor 1. Do not silently floor small
positive scales; nonfinite transformed values also cause explicit failure.
Successful scaling does not waive OLS identification checks. A feature
constant across donors at a required calibration date is confounded with
that date's intercept in time-varying calibration, causing calibration
failure. Constant calibration may identify its common slope from other
calibration dates, subject to full pooled rank after date intercepts.
Uniform, supplied-weight and outcome-only specifications do not acquire
an OLS rank requirement merely because this scaling fallback is used.
This convention is implemented in the new internal helper; R validation
remains pending.

\paragraph{Future-outcome calibration sample.}
Index the regular dates of a training prefix by $t=1,\ldots,T_c$;
$t+1$ means the next declared period, not necessarily the next calendar
integer. Let $Z_{j,t}$ be donor $j$'s transformed feature vector. Relevance
calibration uses only donor pairs $(Z_{j,t},Y_{j,t+1})$ with
$t+1\leq T_c$. Features may include current or earlier outcome observations
when explicitly selected, but not the subsequent response itself.
The target contributes neither features nor responses to calibration;
its selected features enter subsequent matching. No response beyond the
training cutoff may enter calibration, even if it is an untreated donor
outcome. Earlier-date relevance may use the next response inside the
prefix: this is predetermination before donor fitting, not real-time
availability at every earlier date.

OLS without ridge is the agreed calibration estimator. The constant mode
uses a common slope vector and date-specific intercepts:
\[
 Y_{j,t+1}=a_t+\beta^\top Z_{j,t}+\varepsilon_{j,t+1},
 \qquad j\in\mathcal D,\quad t=1,\ldots,T_c-1.
\]
The intercepts enter calibration only, not predictor weights or
counterfactual construction. The pooled slope vector must be uniquely
identified after accounting for these date intercepts. The time-varying
mode uses separate cross-sectional regressions,
\[
 Y_{j,t+1}=a_t+\beta_t^\top Z_{j,t}+\varepsilon_{j,t+1},
 \qquad j\in\mathcal D,\quad t=1,\ldots,T_c-1.
\]
It introduces no rolling windows, time-trend restrictions or coefficient
penalties. Each regression must identify its slopes uniquely; with an
intercept, $J\geq K+1$ is necessary but does not ensure full rank.

\paragraph{Coefficient conversion and the final training date.}
Absolute normalization is the default conversion within empirical
calibration, and squared normalization is the only alternative:
\[
 v_{k,t}^{\rm abs}=\frac{|\widehat\beta_{k,t}|}
 {\sum_{\ell=1}^K|\widehat\beta_{\ell,t}|},\qquad
 v_{k,t}^{\rm sq}=\frac{\widehat\beta_{k,t}^2}
 {\sum_{\ell=1}^K\widehat\beta_{\ell,t}^2}.
\]
Intercepts are excluded. For constant calibration the common converted
vector is repeated across training dates. Uniform weighting remains the
overall fitting default; choosing absolute empirical conversion does not
replace it. Neither conversion is asserted to dominate the other in
counterfactual accuracy. LMG and other relevance allocations are outside
the agreed implementation scope.

In the time-varying mode, the final regression available at this cutoff
uses $Z_{j,T_c-1}$ and $Y_{j,T_c}$. Retain its converted metric for the
last training date:
\[
 V_{T_c}=V_{T_c-1}.
\]
This is an explicit boundary convention in the selected scaled feature
coordinates, not evidence that relevance is constant at the boundary.
The final training observations are retained and $w_{T_c}$ is not forced
to equal $w_{T_c-1}$. Apply this convention separately within every
validation prefix and the final full-pre-period fit. It is distinct from
post-treatment endpoint retention of donor weights.

\paragraph{Failures and remaining scope decisions.}
All-zero slopes or non-unique OLS slopes cause explicit calibration
failure with diagnostics, not automatic uniform weights, predictor
deletion, an arbitrary coefficient selection or ridge. Individual
uniquely estimated zero slopes receive zero weight; other slopes are
normalized without significance filtering or an arbitrary coefficient-size
threshold. The numerical rank rule above is approved; its implementation
and the remaining finite-arithmetic checks require validation.
A failed required calibration fold is reported, not
omitted or shifted. No replacement by an earlier successful metric is
allowed when the final available calibration regression itself fails.

\paragraph{Approved outcome-only specification.}
The outcome-only option uses one contemporaneous outcome feature at each
training date, without covariates or an additional lag stack. Its raw
blocks are $X_{1,t}=Y_{1,t}$ and
$X_{0,t}=(Y_{2,t},\ldots,Y_{J+1,t})$, so $K=1$ and $V_t=[1]$.
All designated training dates enter matching, each with its own donor
vector; the temporal penalty connects these vectors across dates.
This explicitly selected specification bypasses OLS relevance calibration
by design, not as a fallback after calibration failure.
\texttt{tvsc.vmetrics()} still applies the selected scaling and returns
the one-feature metric. Scaling options, the \texttt{sd} default, both
penalty orders, chronological tuning and continuation rules are retained.
Scaling changes the matching--regularization balance at fixed $\lambda$;
reporting outcomes remain in original units. Within validation, only the
training prefix supplies fitting inputs, and subsequent target outcomes
are reserved for scoring. Preparation and metric construction now reflect
this design in source; fitting integration remains unimplemented.
Numerical settings and the empirical comparison layout still need to be
reconciled with these decisions before the complete procedure is implemented.

\subsection{Existing drafts and retained workflow specification}
The source-level contracts and test chronology below document the earlier
drafts. Where their public responsibilities or target-inclusive scaling
differ from Section~\ref{sec:revised-design}, the revised design takes
precedence for future implementation. Earlier passing checks do not
validate the revised functions.

The package defaults are first-difference regularization ($q=1$), uniform
predictor weights ($v_{k,t}=1/K$), chronological pre-treatment selection
of $\lambda$, and projected increment continuation (D). A researcher may
override these choices, including supplying a fixed $\lambda\in[0,1)$.
The agreed scaling default is \texttt{sd}, defined below as a common
scale obtained from within-date variances, not pooled raw levels.
These are software defaults, not a claim that the first order is superior.
Both orders retain equal standing in the empirical comparison.
The public design separates preparation from optional split construction.
For chronological selection, the intended sequence is
\begin{verbatim}
prepared <- tvsc.dataprep(data, ...)
splits <- tvsc.split(prepared, initial = ..., horizon = ..., step = ...)
fit <- tvsc.fit(prepared, q = 1, predictor_weights = "uniform",
                lambda = "cv", splits = splits,
                lambda_grid = seq(0, 0.99, by = 0.01),
                cv_continuation = "projected_increment")
effects <- tvsc.att(fit, prepared,
                    continuation = "projected_increment")
\end{verbatim}
These calls specify the intended complete interface. On September 15,
2026, the five preparation and metric functions were renamed from the
\texttt{panel.*} prefix to \texttt{tvsc.*}, with corresponding source and
test filenames. Arguments, formulas, defaults, return classes and schema
versions are unchanged; no compatibility aliases are provided. Current
names are used below, but the earlier passing transcripts refer to the
former names, not runs of the renamed files. A subsequent September 15,
2026 transcript includes renamed definitions and fixtures: preparation,
splits, blocks and scaling pass. Metrics stops at an incorrect fixture
expectation: removing the scaling field triggers the earlier recipe
check. The fixture now distinguishes an absent field from a null-valued
field; the function body is unchanged. The user subsequently reports a
metrics pass, supplying only the final assertion and success message.
The full rerun transcript, file-loading commands, revisions and session information are
absent; a clean start and independent reproduction are not established.
An initial
raw preparation implementation has been written in \texttt{R/tvsc.dataprep.R};
the user supplied an R console transcript on September 15, 2026 showing the
raw-preparation checks passing. The R version and sourced-file revision were
not supplied, and an independent local rerun remains pending. A separate
split constructor is now written in \texttt{R/tvsc.split.R}, with R
fixtures in \texttt{tests/test-tvsc.split.R}. A second user-supplied R
transcript on September 15, 2026 shows these split checks passing, including
the eight-date example and input-validation cases. Source revisions and
R session information remain unrecorded; neither run was independently
reproduced in the manuscript workspace. A raw block builder is now written
in \texttt{R/tvsc.blocks.R}, with fixtures in
\texttt{tests/test-tvsc.blocks.R}. A third user-supplied R transcript on
September 15, 2026 shows the block checks passing, including preservation
of pre-scaled predictors and transformed-outcome passthrough. The transcript
does not include source-loading commands, initial panel creation, source
revisions or R session information; a clean-session reproduction and the
standalone example remain pending. A separate scaling draft is now written
in \texttt{R/tvsc.scale.R}, with fixtures in
\texttt{tests/test-tvsc.scale.R}. A fourth user-supplied R transcript on
September 15, 2026 shows the displayed scaling checks passing, including
all four methods, zero-scale rules, cutoff isolation and numerical failure
cases. Fixture creation is included, but source-loading commands, exact
revisions and R session information are absent. Independent clean-session
reproduction and the standalone example remain pending.
The predictor-metric constructor is now written in
\texttt{R/tvsc.metrics.R}, with fixtures in
\texttt{tests/test-tvsc.metrics.R}. It constructs uniform weights or validates
supplied weights; it does not perform empirical calibration. The reported
corrected-suite pass is supported only by the supplied ending, not a full
rerun transcript. Independent reproduction and the standalone example
remain pending.
The revised calibration source is described above and awaits R execution;
fitting, prediction and effect functions remain unimplemented.
A solver-independent primary-problem builder is now written in
\texttt{tvsc-problem.R.txt}, with fixtures in
\texttt{tvsc-problem-test.R.txt}. It constructs dense reference matrices
and constraints for both orders, not fitted weights. R execution of this
new suite and its standalone example remains pending; the earlier
transcripts do not validate it. This builder is set aside while the revised
preparation and calibration contracts are settled; it is not the immediate
development checkpoint. The five-function GitHub handoff has not
yet produced integration results in this workspace.
A prediction method separates counterfactual construction from
the treatment-effect calculation.
If the researcher supplies a fixed penalty, no validation split is required:
\begin{verbatim}
fit <- tvsc.fit(prepared, lambda = 0.4)
\end{verbatim}
The value 0.4 illustrates a fixed input, not a recommended penalty. This
call uses the full designated pre-treatment sample. It yields training
diagnostics, not an out-of-sample assessment.

\subsection{From panel data to fitted donor paths}
The revised raw preparation draft requires explicit predictor names,
an integer-valued numeric period index and a declared regular step.
It returns raw training data, schema metadata and separately stored future
donor and target outcomes, plus aligned full-pre-period raw blocks. It does
not accept or estimate a scaling recipe or construct splits. Constant raw predictors
are retained and flagged. These implementation limits are recorded in the
strategy note; the paragraphs below describe the intended complete contract.

Begin with a longitudinal table containing unit, date, outcome and predictor
columns. Specify the target, donor eligibility, intervention date and
eligible pre-treatment period. The intervention boundary restricts which
target observations may enter fitting; it is not a training/validation
split within that period. Preparation requires no validation windows or
reporting horizon. Validate unique unit--date records, common
donor ordering, regular spacing and the availability of all required cells.
The initial workflow rejects unresolved missingness rather than silently
changing the panel. The proposed \texttt{tvsc.dataprep()} retains raw
pre-treatment observations and the block-construction and scaling recipes.
It does not choose cutoffs, fit preprocessing parameters or withhold
observations for assessment. An internal training builder estimates any
data-dependent transformations for the actual fitting sample: the full
designated pre-period for direct fitting, or each prefix during validation.
Any post-treatment
target observations are stored separately and are not exposed to fitting
or tuning. Full-pre-period transformations must not be reused within
shorter validation prefixes.

The separate \texttt{tvsc.split()} function constructs date assignments
and fold metadata from the prepared object. It does not scale predictors,
calibrate metrics, fit weights or return independently transformed datasets.
Each fold records its training dates, cutoff and subsequent validation
dates, selecting all eligible units at each date. Windows must be within
the eligible pre-period, with at least three training dates and a nonempty
validation window strictly after the fitting cutoff. The agreed interface is
\begin{verbatim}
splits <- tvsc.split(prepared, initial, horizon, step = 1)
\end{verbatim}
Training windows expand from the first eligible date. The initial size and
validation horizon are required inputs counting dates, not panel rows;
the stride defaults to one date. The initial implementation retains only
full validation windows on that stride, without shortening the final window
or adding an off-stride cutoff. It reports assessment counts by date and
all unassessed dates. With eight pre-period dates, an initial size of four
and horizon of two yield training/validation positions 1--4/5--6,
1--5/6--7 and 1--6/7--8. Repeated validation dates are not independent
replications; later scoring must specify their weighting. The split object
contains a schema copy for downstream compatibility checks, not a checksum
of the data. It does not reserve an independent test set automatically.
Chronological tuning requires valid splits, whereas a fixed-penalty fit
does not. A CV request without an explicit split specification must request that
information rather than silently create windows or use the training loss.
Splits can also support held-out assessment, but tuning scores are not an
independent evaluation of the tuned procedure.

\paragraph{Raw predictor-block assembly.}
The agreed next layer separates matrix assembly from scaling. Its initial
implementation uses the already declared contemporaneous predictors:
\begin{verbatim}
blocks <- tvsc.blocks(prepared, cutoff = 2003)
\end{verbatim}
The cutoff is an explicit date in the designated pre-period, not a row
number; the prefix must contain at least three dates. For a full fit, supply
the last declared pre-period date. For a validation fold, supply that fold's
cutoff. The builder reuses \texttt{tvsc.dataprep()} to validate and order
only the selected prefix. Schema and stored time indices are checked before
selection; required numerical values and unit--date completeness are checked
within the prefix, not across later validation values.

The returned \texttt{tvsc\_blocks} object stores date-named lists
\texttt{X1} and \texttt{X0}: named length-$K$ target vectors and $K\times J$
donor matrices, respectively. Predictor and donor ordering follow the
declared schema; $K=1$ does not drop donor-matrix dimensions. As-supplied
outcomes are retained separately as a date-named vector \texttt{Y1} and a
date-by-donor matrix \texttt{Y0}, even when the outcome is not a matching
predictor. The object includes a prefix-restricted schema, used dates,
cutoff, dimensions, recipe and prefix-only constant-predictor diagnostics.
It contains no later observed values or future tables. No lags, summaries,
scaling, predictor weights or donor weights are constructed. This raw
builder does not replace the later training-specific transformation layer.

\paragraph{No additional scaling and externally transformed inputs.}
Here ``raw'' means as supplied to preparation, not necessarily never
transformed. Both preparation and block assembly already preserve supplied
values without additional centering or scaling. The agreed future scaling
interface will support \texttt{scaling = "none"}, meaning no additional
scaling whether inputs are in economic units or have already been
transformed. It overrides the agreed \texttt{sd} default described below.
This argument is implemented in the separate scaling draft, not accepted
by the existing preparation or block functions. Their
\texttt{recipe\$scaling = NULL} records no package-applied scaling; it is
not a claim about upstream processing and should not be replaced manually.

External preprocessing must be documented: affected columns, transformation,
parameters common to target and donors, output units, and dates and units
used to estimate the parameters. A learned transformation must use only
the relevant training prefix during validation. Selecting a prefix from
previously full-period-standardized values does not remove that information
leakage. A fixed externally specified unit conversion requires no such
full-period estimation. The current functions neither verify provenance
nor reverse earlier transformations.

The outcome column should remain in the intended reporting units. A
transformed outcome for predictor matching can be supplied separately, for
example \texttt{outcome = "sales"} with predictors
\texttt{sales\_scaled} and \texttt{income\_scaled}. The stored outcome
arrays preserve the supplied outcome column, not automatically recovered
original units. Upstream transformations of that column therefore require
an explicit interpretation or a separately specified back-transformation.

\paragraph{Agreed initial scaling scope.}
This paragraph records the target-inclusive source draft, superseded for
all revised weighting modes by donor-only learned scaling in
Section~\ref{sec:revised-design}; its zero-contrast statements refer to
that original reference population.
The initial scaling draft offers only \texttt{sd} (the default),
\texttt{zscore}, \texttt{none} and \texttt{custom}. The source and fixtures
are written, with a user-reported passing test transcript and independent
reproduction pending. Pooled-level alternatives,
MAD, IQR, min--max, maximum-absolute-value and other robust methods are
deferred. The short names have the following panel-specific definitions.

For a training prefix of length $T_c$, let $N=J+1$ and let $x_{i,k,t}$
denote predictor $k$ for the target and eligible donors. Define
\[
 \bar x_{k,t}=\frac1N\sum_{i=1}^N x_{i,k,t},\qquad
 s_{k,t}^2=\frac1{N-1}\sum_{i=1}^N(x_{i,k,t}-\bar x_{k,t})^2,
 \qquad s_{k,c}=\left(\frac1{T_c}\sum_{t=1}^{T_c}s_{k,t}^2\right)^{1/2}.
\]
The \texttt{sd} option divides every matching value by $s_{k,c}$,
without centering. It uses one common scale per predictor throughout
the prefix: the root mean square of the within-date standard deviations,
not the standard deviation of pooled raw levels. The \texttt{zscore}
option instead maps each value to $(x_{i,k,t}-\bar x_{k,t})/s_{k,t}$,
using a separate center and scale at each training date. The \texttt{none}
option preserves supplied values; \texttt{custom} divides by positive,
finite researcher-supplied scales matched by predictor name and common
across dates. Scaling changes matching blocks only, leaving reporting
outcomes in their supplied units.

All transformations are common across units at a given date and fixed
before donor estimation. Learned parameters use only the actual training
prefix and are reconstructed within each validation fold. No future target
values enter scaling; continuation uses donor outcomes in reporting units,
not future target-based z-scores. Store the method, reference units, dates,
centers, scales and diagnostic decisions.

The agreed call is
\begin{verbatim}
scaled <- tvsc.scale(blocks, scaling = "sd", scales = NULL)
\end{verbatim}
It accepts raw blocks and returns a new object inheriting from
\texttt{tvsc\_scaled\_blocks} and \texttt{tvsc\_blocks}. A non-null
\texttt{recipe\$scaling} records the method, formula identifier, reference
units, fitting dates, cutoff, applied centers and divisors, estimated
within-date SDs when applicable, and fallback diagnostics. Parameter
matrices have predictor rows and date columns. Custom scales must be named
exactly once for every predictor and are aligned by name; they are rejected
for other methods. Every output, including \texttt{none}, is marked as
processed and cannot be passed through scaling again. Another choice must
start from raw blocks. The identity method records zero centers and unit
divisors without arithmetic on the input predictor arrays.

Under \texttt{sd}, exact equality across all reference units at every
training date permits a divisor-1 fallback, recorded for all dates of that
predictor. Under \texttt{zscore}, exact equality at an individual date uses
its common value as center and divisor 1, producing zeros. These cells
have no matching contrast; predictors are retained and predictor weights
are not renormalized. A zero or nonfinite computed SD despite unequal
values is a numerical error, not a fallback. Positive scales are neither
clipped nor floored; nonfinite transformed values are rejected. The
implementation records the smallest positive applied divisor and fallback
count, without certifying solver conditioning. Original raw-block
diagnostics and reporting outcomes are preserved unchanged.

Common centering cancels from simplex-weighted discrepancies, but scaling
does not. If $e_{k,t}$ is a supplied-unit discrepancy, for positive
estimated scales the effective
coefficient is $v_{k,t}/s_{k,c}^2$ under \texttt{sd} and
$v_{k,t}/s_{k,t}^2$ under \texttt{zscore}; fallback cells instead use the
applied unit divisor. The latter method therefore changes
date-specific matching priorities; these options are not merely centered
and uncentered versions of the same criterion. Uniform predictor weights
mean equal coefficients in the chosen scaled units, not equal economic
importance. Within-date dispersion excludes common time shifts that cancel
from matching discrepancies but can inflate pooled-level dispersion.
This motivates the default without establishing optimality or statistical
consistency. The conditional theory applies to the resulting finite
transformed blocks; scaling alone supplies no asymptotic guarantee.

Next construct predictor blocks and transformations within the designated
training sample. Supply uniform predictor weights or run the separately
specified empirical calibration, then freeze $\mathcal V$. A calibration
must define its response, regressors, dates, pooling or time variation,
and rules for rank deficiency and zero normalization denominators. Its
response must not also appear as an identical predictor. A common pooled
regression does not by itself specify date-varying coefficient vectors.
For an empirical metric, retain the calibration output and its provenance
so that it can be reproduced independently of donor fitting.

The initial \texttt{tvsc.metrics()} interface separates this validation
step from empirical calibration. Let $T_c$ denote the number of dates in
the supplied training prefix. The three accepted input shapes produce
\[
 v_{k,t}=\begin{cases}
  1/K, & \text{uniform},\\
  a_k, & \text{supplied common vector }a,\\
  M_{k,t}, & \text{supplied }K\times T_c\text{ matrix }M.
 \end{cases}
\]
The vector and matrix must have finite nonnegative entries and unit column
sums, subject to an explicitly recorded numerical tolerance. Vector names
must cover the predictors; matrix row and column names must cover the
predictors and prefix dates exactly once. Order may differ; missing or
extra labels are rejected. The source draft accepts an absolute sum error
of at most \texttt{tolerance = 1e-8} by default and requires positive sums.
It preserves accepted values without renormalization or clipping, so
numerical acceptance does not assert exact unit sums. Zero entries are
allowed. These are input shapes, not automatic labels of empirical origin.
\begin{verbatim}
scaled <- tvsc.scale(tvsc.blocks(prepared, cutoff), "sd")
metrics <- tvsc.metrics(scaled, predictor_weights = "uniform")
\end{verbatim}
Use \texttt{tvsc.scale(..., "none")} to declare no scaling explicitly.
The returned \texttt{tvsc\_metrics} object stores a canonical
predictor-by-date weight matrix, schema, cutoff, scaling record and
diagnostics. It reads no numerical predictor or outcome values to construct
weights and does not validate those arrays. Alignment checks do not
authenticate the scaling record or establish absence of leakage from
upstream calibration. An optional \texttt{provenance} description is stored
as unverified. The eventual fitter must validate numerical arrays and
compatibility with these metrics. Supplied empirical weights must still
be recalibrated within each fitting prefix; slicing a full-period
calibration is not sufficient. Both orders and all candidate penalties
share the same frozen prefix-specific blocks and metrics.

For each $q$ and candidate $\lambda$, construct the matrices in
Section~\ref{sec:framework} and solve the quadratic program
\[
 \min_{\theta\in\mathcal C}\frac12\theta^\top P_q\theta+c^\top\theta,
\]
\[
 P_q=2\{(1-\lambda)\mathcal A^\top\mathcal A+
                      \lambda\mathcal D_q^\top\mathcal D_q\},
 \qquad c=-2(1-\lambda)\mathcal A^\top b.
\]
The omitted term $(1-\lambda)b^\top b$ is constant. The source draft
\texttt{tvsc.problem(blocks, metrics, lambda, q = 1)} retains it explicitly,
returning \texttt{P}, \texttt{linear} and \texttt{constant} so that
\[
 F(\theta)=\tfrac12\theta^\top P_q\theta+c^\top\theta
                         +(1-\lambda)b^\top b.
\]
It also retains $\mathcal A$, $b$ and $\mathcal D_q$. Coordinates stack
donors in their declared order within each date. With
$S=I_T\otimes\mathbf1_J^\top$, its constraints are
$S\theta=\mathbf1_T$ and $0\leq\theta\leq1$; the upper bounds are
redundant but explicit. Differences count declared period steps, not
calendar-index distances. Metrics and scaling metadata must agree, and
the builder revalidates numerical predictor arrays. It does not read or
numerically validate reporting outcomes, calibrate metrics or solve either
selection problem. Dense storage is intended for small reference checks,
not a measured performance claim. Nonfinite arithmetic produces an error.
The stored $\mathcal D_q$ is unweighted even at $\lambda=0$; a later
secondary solve must use its $\sqrt\lambda$ scaling as specified in the
appendix, not impose an unscaled difference-image equality.

Numerical singularity is not a
reason to add an unreported ridge term. Apply the secondary selection
rule, with explicit tolerances for feasibility, objective accuracy and
preservation of the primary solution images; its exact formulation is
given in Supplement~\ref{app:selection}. Retain failure statuses rather
than replacing unsuccessful fits with a different estimator.

By default, select $\lambda$ by chronological pre-treatment validation.
At each cutoff, repeat block construction, scaling and any
predictor-weight calibration using that prefix alone, fit the donor path,
and predict the subsequent untreated observations using D. For a fixed
order and predictor-weighting specification, let $c_f$ denote a fitting
cutoff and $\mathcal H_f$ its subsequent untreated validation dates. Use
\begin{equation}\label{eq:lambda-validation}
 \widehat\lambda\in\arg\min_{\lambda\in\mathcal G}
 \frac1{|\mathcal F|}\sum_{f\in\mathcal F}\frac1{|\mathcal H_f|}
 \sum_{t\in\mathcal H_f}
 \left(y_t-x_t^\top\widehat w^D_{t\mid c_f}(\lambda)\right)^2.
\end{equation}
The weights in this expression are refitted within each prefix using its
predetermined metrics. Target validation outcomes enter only the score;
contemporaneous donor outcomes enter prediction. Select $\lambda$ separately
for each dynamic specification, using common validation dates and horizons.

The default grid is $\mathcal G=\{0,0.01,\ldots,0.99\}$, containing
100 values and shared across the two orders. Its equivalent additive
coefficients range from zero to 99 and are not equally spaced. This is
a finite search range, not the entire theoretical domain. Other declared
finite grids in $[0,1)$ are permitted. Record any upper-bound selection;
endpoints, spacing, the tie rule and any boundary-triggered expansion must
be specified before assessment and use training information only.
A candidate is ineligible if any required validation fit fails its
acceptance checks; do not average only its successful folds. If none is
eligible, report selection failure rather than changing validation dates.
Predictor-weighting modes remain declared specifications, not variables
jointly optimized with donor weights. If a researcher instead chooses E
as the prediction procedure, its tuning score must use E consistently.

Scores used for tuning are validation scores, not independent tests.
Where the pre-period allows, reserve additional assessment dates outside
tuning. After configuration choices are fixed, repeat the entire workflow
on the full designated pre-treatment period. Store the predetermined metrics,
selected penalty, donor paths and diagnostics. Within-fit E/D comparisons
use the same fitted configuration; separately tuned E/D comparisons must
be labeled as comparisons of complete procedures.

The proposed \texttt{tvsc.fit()} returns a model object rather than only
a weight matrix. It stores donor identities and ordering, the fitted path,
endpoint and final increment, predetermined metrics and transformations,
$q$ and selected $\lambda$, training predictions and residuals, RMSE and
MAE, fold-level validation scores, cutoffs and horizons, and numerical
acceptance diagnostics. The stored \texttt{cv\_continuation} identifies
the extension used during tuning. Empirical calibration remains a separate,
explicitly supplied procedure, not an implicit optimization over predictor
weights inside the donor solver.

\subsection{From fitted paths to treatment effects}
Extend each accepted donor path using \eqref{eq:continuation}, and form
untreated predictions from contemporaneous donor outcomes. A static fit
retains its common donor vector. A rejected dynamic fit makes both of its
continuation outputs unavailable. No post-treatment target observation
can revise the metrics, regularization or donor weights.

The proposed prediction method accepts the fitted object, aligned donor
outcomes as \texttt{newdata}, and a continuation rule. It requires no target
outcomes. The
\texttt{tvsc.att()} function uses this same prediction operation and then
joins the observed target outcomes. Its continuation argument accepts
\texttt{projected\_increment} (D, the default) or \texttt{endpoint} (E),
without refitting or recalibrating. Results record both the continuation
actually used and the fitted object's tuning continuation. Choosing E for
a D-tuned object is an endpoint prediction from a D-tuned fit, not an
independently tuned E procedure. The latter requires a new fit with
\texttt{cv\_continuation = "endpoint"}.

Only after the predictions are fixed, join the observed treated outcomes
and compute
\begin{equation}\label{eq:effects}
 \widehat\tau^R_{T+h}=Y^{\rm obs}_{1,T+h}
                         -\widehat Y^R_{1,T+h}(0),\qquad
 \widehat{\operatorname{ATT}}^R_H
                 =\frac1H\sum_{h=1}^H\widehat\tau^R_{T+h}.
\end{equation}
The corresponding target is
$\operatorname{ATT}_H=H^{-1}\sum_{h=1}^H\tau_{T+h}$.\footnote{For this
single-treated-unit design, ATT denotes an average over a declared set
of post-treatment dates, not a population average across treated units.}
All methods use the same declared complete reporting window. Missing
predictions or target outcomes make that window's estimate unavailable
unless a separate missing-data procedure has been specified.

When the observed target outcome equals $Y(1)$,
\[
 \widehat{\operatorname{ATT}}^R_H-\operatorname{ATT}_H
 =\frac1H\sum_{h=1}^H
       \{Y_{1,T+h}(0)-\widehat Y^R_{1,T+h}(0)\}.
\]
Consequently, valid bounds on absolute counterfactual errors also bound the
absolute ATT error by their average. This is a conditional implication,
not an estimated standard error. Interpreting the gaps causally requires
credible untreated donors, no relevant spillovers or anticipation, and
an adequate untreated relationship across the cutoff \cite{abadie2021}.

The package should return fitted and extended donor weights, supplied
predictor weights, outcome paths, annual gaps, average effects, fit and
validation scores, and numerical diagnostics. Plotting functions should
use common axes and reporting windows across methods. Data preparation,
calibration, fitting, prediction and scoring remain separate operations,
with saved settings sufficient to reconstruct each result.


\section{Implementation and validation to be completed}
The first implementation will need to validate preparation without splits,
the separate date-based split constructor and information separation,
training-specific preprocessing, both penalty orders, chronological tuning,
secondary selection and both continuation rules. Deterministic test
fixtures should check simplex feasibility, objective and matrix scaling,
known small-problem solutions, failure handling and agreement of effect
calculations with saved predictions. These software checks are distinct
from a Monte Carlo performance study; no such study is proposed here.

Timing and numerical acceptance results must be measured after the
implementation exists, not inferred from convexity alone. Reproducible
examples and environment records will accompany the software paper.
The substantive California results, eight-case comparison and treatment
effect interpretation belong in the methodological paper, using the
validated functions and a recorded software version.

\section{Conclusion}
The workflow separates operations that use different information and
serve different purposes. Predictor calibration precedes donor fitting;
validation uses untreated assessment outcomes only for scoring; and
counterfactual prediction precedes treatment-effect calculation. This
paper will document how the implementation enforces these distinctions.
Preparation is usable without validation, while chronological penalty
selection requires declared splits. The current focus is to settle these
remaining fitting contracts and validate the staged implementation. The
revised interface combines raw preparation and block construction, and
combines scaling and predictor metrics under \texttt{tvsc.vmetrics()}.
Donor-only future-outcome OLS calibration, absolute conversion by default,
date-by-date time-varying slopes and final-date metric retention are now
specified as design choices. Preparation/block integration and split-input
compatibility are now revised in source, with both fixture suites passing
in the supplied user-run R transcript. Independent reproduction remains
pending. The new metric-construction source now implements donor-only
scaling and calibration, with a user-reported fixture pass supported only
by the final isolated-loading excerpt; the full log and independent
reproduction remain pending. Fitting,
prediction and effect operations remain unimplemented; earlier user-run
checks concern the previous source drafts. The
methodological paper remains the theory baseline.

\begin{thebibliography}{9}
\bibitem{abadie2021}
Abadie, A. (2021). Using synthetic controls: Feasibility, data requirements,
and methodological aspects. \emph{Journal of Economic Literature},
59(2), 391--425. DOI: \nolinkurl{10.1257/jel.20191450}.
\end{thebibliography}

\appendix
\section{Numerical selection contract}\label{app:selection}
Given an exact primary minimizer $\theta_p$, the exact secondary problem is
\[
 \min_{\theta\in\mathcal C}\|\theta\|_2^2
 \quad\text{subject to}\quad
 \mathcal A\theta=\mathcal A\theta_p,\qquad
 \sqrt\lambda\mathcal D_q\theta=\sqrt\lambda\mathcal D_q\theta_p.
\]
The methodological paper proves that these image constraints characterize
the primary minimizer set. At $\lambda=0$ the difference-image condition is
vacuous; imposing the unscaled condition would change the estimator.
Approximate image preservation must not be described as certification of
exact minimum-norm selection. Acceptance tolerances and failures must be
recorded. Because the mixing objective is $1-\lambda$ times its equivalent
additive objective, objective-gap tolerances must be rescaled when comparing
those representations. The primary and secondary numerical algorithms
and their validation remain implementation work.
\end{document}